\begin{tikzpicture}[font=\small,line cap=round,line join=round]\draw[otOrange,dashed,line width=1pt](3.000,-0.300)--(3.000,2.800);\draw[otNavy,line width=1.2pt](1.000,0.000)--(1.000,2.400);\fill[otNavy!30](1.000,2.400)--(2.400,2.000)--(1.000,1.600)--cycle;\draw[geoTrait](1.000,2.400)--(2.400,2.000)--(1.000,1.600)--cycle;\draw[otNavy,line width=1.2pt](5.000,0.000)--(5.000,2.400);\fill[otNavy!30](5.000,2.400)--(3.600,2.000)--(5.000,1.600)--cycle;\draw[geoTrait](5.000,2.400)--(3.600,2.000)--(5.000,1.600)--cycle;\node[below,font=\tiny,inner sep=1.2pt,text=otInk]at(1.000,-0.050){drapeau};\node[below,font=\tiny,inner sep=1.2pt,text=otInk]at(5.000,-0.050){image};\end{tikzpicture}\par\vspace{1.5mm}{\otRegular\fontsize{8}{10}\selectfont\color{otMuted}Un drapeau et son symétrique par rapport à l’axe en tirets.\par}
